% =====================================================================
% BRANCH B DRAFT: Section 3 - quality-constrained comparison
% Working title: When Debiasing Scores Reward Broken Text
% Target: ARR October 2026 cycle (submission Oct 12)
%
% DRAFTING RULE FOR THIS FILE: no hedging in the body. Every scope
% condition is stated once, plainly, and the rest live in Section 6
% (Limitations). Numbers are from analysis-output/sweep_points.csv and
% analysis-output/stats-appendix.md.
% =====================================================================

\section{What the methods do once quality is held constant}
\label{sec:frontier}

Section~\ref{sec:audit} showed that a high target-metric score can be produced by
damaging the response. That makes the published ranking uninterpretable on its own:
it does not distinguish an intervention that removed the bias from one that removed
the answer. This section re-runs the comparison with response quality measured
alongside the target metric.

\paragraph{Protocol.} This sweep is exploratory; Section~\ref{sec:confirm} re-tests its
claims under a pre-registered design. We sweep the steering strength of every method across its full
published range and generate 250 AccessEval responses at each setting, giving 87 sweep
points and 21{,}750 responses in total. Each response receives a medicalization score
and a quality score from an LLM judge (Llama-3.1-8B-Instruct, temperature $0.1$, seed
42) on a ten-point accessibility rubric. Ten of the 87 points apply no steering; these
give a reference of $7.92 \pm 0.04$ for quality and $0.659 \pm 0.043$ for
medicalization, where the intervals are standard deviations across the ten points. A
medicalization score of zero is neutral, and scores below zero mean the model withholds
medical information from users who asked for it.

\paragraph{Flagging degenerate output.} We mark a sweep point as degenerate from
properties of the text and the target metric alone, never from the quality score: fewer
than half the 250 responses distinct, more than $10\%$ empty or unparseable, or a mean
medicalization below $-0.05$. Keeping the quality score out of the definition means the
judge remains an independent outcome rather than a restatement of the flag. 38 of the
87 points are flagged, and the analysis below is insensitive to where the three cut-offs
sit: across all 27 combinations of $\{0.3, 0.5, 0.7\}$, $\{0.05, 0.10, 0.20\}$ and
$\{-0.10, -0.05, 0\}$ the Pareto front holds the same twelve configurations.

\subsection{The cost at each published operating point}

At the configuration each method's authors recommend, all five reduce medicalization and
all five cost quality, but the cost spans an order of magnitude. SADI posts the lowest
medicalization score of the five and loses $60\%$ of the model's judged quality. Angular
reaches $-0.222$, past neutral, with token corruption inside words. Ordering the methods
by the target metric and by quality gives different rankings.
Table~\ref{tab:e1} gives the pre-registered numbers for these points, and
Appendix~\ref{app:earlier} the earlier single-setting evaluation with sentiment and
regard.

\subsection{The frontier}

Figure~\ref{fig:frontier} plots all 87 points. The intact points trace a shallow curve
that stays near the unsteered reference until roughly $0.2$ on the target metric, then
falls away. Every point that reaches or passes neutral is degenerate.

\begin{figure}[t]
\centering
\includegraphics[width=\linewidth]{figure-01-quality-debiasing-frontier.pdf}
\caption{Judged quality against medicalization for 87 sweep points, $n{=}250$ responses
each. Filled markers are intact output, hollow markers are degenerate by the criterion
in Section~\ref{sec:frontier}, which uses text properties and not the quality score.
Error bars are bootstrap $95\%$ confidence intervals on the mean quality. The grey band
is the unsteered reference, $\pm 2$ standard deviations across ten no-steer points. The
shaded region past neutral marks responses that withhold medical information, which is
not an improvement. The blue line is the Pareto front over intact points.}
\label{fig:frontier}
\end{figure}

The Pareto front over intact points holds twelve configurations, ten of them CAD
operators. Restricting to points that remove at least half the baseline medicalization,
all four surviving front configurations are CAD: the probe at $\alpha{=}9$, the
mean-difference operator at $\alpha{=}8$ and $\alpha{=}7$, and the projection at
$\alpha{=}1$. No baseline method reaches that region of the plot without degenerating.

At matched strength the projection also scored above CAA in this sweep, by $0.30$
$[0.07, 0.53]$. That difference does not survive the held-out test
(Section~\ref{sec:confirm}), and we do not rely on it.

\subsection{Collapse is a threshold, and the published settings sit close to it}

Figure~\ref{fig:collapse}(a) in the appendix sweeps steering strength for each method. Quality holds
close to the unsteered reference and then falls steeply over a narrow range of
strengths. The fall is not gradual in any of them.


Raw strengths do not compare across methods, since each parameterises its intervention
differently. What does compare is the margin between the setting an author recommends
and the setting at which the method breaks. Taking the breaking point as the first
strength whose mean quality falls below $6.0$, that margin is $3.0\times$ for the CAD
projection, $2.0\times$ for CAA and for FairSteer, $1.9\times$ for the CAD
mean-difference operator, $1.2\times$ for the CAD probe, and $1.0\times$ for SADI.

SADI's published setting is its breaking point. At $s{=}5$ it scores $7.5$ on quality
and removes little bias; at $s{=}10$, the value its authors report, quality is $3.2$.
There is no intermediate setting in the published sweep at which the method both works
and leaves the text intact.

A wider margin is not immunity. The CAD projection is flagged degenerate from
$\alpha{=}3$, where $70\%$ of its responses become unparseable, and at $\alpha{=}5$ it
returns 75 empty responses out of 250. Every method here has a range in which it works
and a point past which it destroys the output. What separates them is how much room
sits between that point and the strength needed to remove the bias.

This also explains why the leaderboard inverts the ranking. Methods are tuned to
maximise the target metric, and the target metric keeps rewarding larger steering
strengths past the point where the text breaks, because broken text contains neither
medical nor professional vocabulary and the log-odds contrast
\citep{monroe2008fightin} goes to zero. The degeneration modes themselves are familiar
from open-ended generation \citep{welleck-etal-2019-neural}; what is new here is that a
fairness metric scores them as success.

\subsection{A second architecture}

On Qwen2.5-7B the projection operator reaches
$d{=}0.392$ at $\alpha{=}1$ with no degenerate output, and collapses at $\alpha{=}2$
rather than the $\alpha{=}3$ we observe for Llama-3.1-8B. The operating window is
narrower on Qwen, and the threshold structure of Figure~\ref{fig:collapse} replicates
with a different threshold.

\paragraph{Scope.} The sweep draws a separate 250-prompt subset per method and fits every
method on a pool that contains its evaluation prompts, so it supports the shape of the
curve and the ordering rather than paired effect sizes. Section~\ref{sec:confirm} removes
both limitations.
